% Introduction of the two-readings note.
\section{Introduction}
\label{r0:sec:main}

The letter $X$ is used in the arena of \cite{ES17} for two different things. In the momentum--force construction it is a displacement,
a vector of $V=\R^{3,3}$ with the dimension of a length; in the field-strength reading of Section~12 of \cite{ES17} the same vector space
carries the electric and magnetic fields, $X=\sum_i(B_ie_i+E_ie_{i+3})$. A working note of 24 October 2026 \cite{FMB} proposed to keep the two
readings strictly apart, $X_{\rm cl}$ for the displacement and $X_{\rm EM}=(\bB,\bE)$ for the six components of the field two-form, and to
build the four-momentum of each separately. This note answers four questions put on that basis:
\begin{enumerate}[label=(\arabic*),leftmargin=2.2em]
\item construct the Lagrangians for $X$, for the Dirac field and for the energies;
\item differentiate between the two $X$'s, the first a length and the second the pair $(\bB,\bE)$;
\item construct the four-momentum, the length-valued $X$ being taken along the three directions of the spatial triple $V^-$, of signature
$(-,-,-)$;
\item place $X_{\rm EM}=(\bB,\bE)$ on the octants, each of the six components $E_x,E_y,E_z$, $B_x,B_y,B_z$ on an octant.
\end{enumerate}
Question (4) is read as the pair: $X_{\rm EM}$ is not a sum of two vectors but the six independent components of one field strength.
Questions (1) and (2) are answered in \cref{r1:sec:main}, question (3) in \cref{r2:sec:main} and question (4) in \cref{r3:sec:main}.
\Cref{r4:sec:main} takes up the element that the answers share: the temporal volume $T=e_1e_2e_3$, which is the Dirac mass form of the
classical reading for every clock $n$, through $T=nJ(n)$ with $J(n)$ the Dirac complex structure of that clock, and the duality operator
of the field reading on the octants. \Cref{r5:sec:main} collects the answers.

\paragraph{What is established, in brief.} The two readings share a vector space and a Clifford action, but their units, their natural
Lorentz algebras, their canonical conjugates and their observables differ, and no equivariant linear map relates them
(\cref{r1:thm:type}). The Lagrangians are the six-dimensional worldline and the interior shifted Lagrangian for $X_{\rm cl}$, the real
Grassmann Lagrangian for the Dirac field, and the Maxwell Lagrangian $-\tfrac12\Re q$ for the field, whose Hamiltonian is the undivided
energy $E_6$; the rest energy $E_0=\tfrac12|q|$ is a Lorentz-invariant and conformal Lagrangian, but not a duality-invariant one and not an
admissible Hamiltonian (\cref{r1:thm:L0,r1:prop:ham}). With a clock $n$ on the temporal triple, the four-momentum $(E/c,\bp)$ has $\bp$
conjugate to the three lengths of $V^-$ and $E/c$ conjugate to the time read by $n$, and the four-dimensional mass read by $n$ never exceeds
the six-dimensional one (\cref{r2:prop:cons,r2:prop:rotate}); the field's four-momentum comes from its stress tensor and is a four-vector
only for an isolated whole (\cref{r2:sec:field}). On the octants, the joint eigenlines of the pole operators, $\bE$ sits on the three
octants of Fock weight one and $\bB$ on the three of weight two; Maxwell's four equations become one first-order equation with the two
Gauss laws on the two poles; and the boosts and the duality rotation of the field act on the octants only through odd operators, so that
the octant representation of the field is not the spin representation in which a Dirac field transforms (\cref{r3:thm:noeven,r3:prop:odd}).

\paragraph{Conventions and classes.} The conventions are those of \cite{ES17}: $\eta=\operatorname{diag}(1,1,1,-1,-1,-1)$, the temporal
triple $V^+=\langle e_1,e_2,e_3\rangle$ carries the converted times of the classical reading and $\bB$ of the field reading, the spatial
triple $V^-=\langle e_4,e_5,e_6\rangle$ the lengths and $\bE$; the paper's $\bB$ is minus the textbook magnetic field. Each section states
its units. Every claim carries a class stated in words: \emph{fixed} (an identity or theorem that could have failed), \emph{constructed}
(admitted and exhibited, not selected), \emph{obstructed} (impossible under stated hypotheses), \emph{open}, \emph{inherited} (true of
ordinary physics in any notation, cited), \emph{interpretive} (a reading), or \emph{restated} (established in \cite{ES17}, \cite{RE} or
\cite{FMB} and re-verified here). Where \cite{FMB} states something, the text says whether it is confirmed, corrected or refuted.
Every finite-dimensional and symbolic claim is checked by a script; the checks are described in \cref{r9:sec:main}.
